\documentclass[A4,10pt,twoside]{book}
\usepackage{amd}
\usepackage[fontset=fandol]{ctex}
\usepackage{amsmath}
\usepackage{bm}
\usepackage{amssymb}
\usepackage{esint}
\usepackage{fontspec}
\usepackage{graphicx}
\usepackage{hyperref}
\usepackage{tasks}
\renewcommand{\theHchapter}{\thepart.\thechapter}
\setmainfont{Times New Roman}
\setsansfont{Arial}
\setmonofont{Courier New}
\setcounter{tocdepth}{0}  % 只显示chapter级别的内容
% 设置所有行内公式为行间公式的样式
\everymath{\displaystyle}
% %--------------------------------------------------------------------------
% %         General Setting
% %--------------------------------------------------------------------------

\graphicspath{{images/}{../images/}} %Path of figures
\setkeys{Gin}{width=0.85\textwidth} %Size of figures
\setlength{\cftbeforechapskip}{3pt} %space between items in toc
\setlength{\parindent}{0.5cm} % Idk
\input{theorems.tex} % Theorems styles and colors
\usepackage[english]{babel} %Language

\setlist[itemize]{itemsep=5pt} % Adjust the length as needed
\setlist[enumerate]{itemsep=5pt} % Adjust the length as needed



% \usepackage{lmodern} %  Latin Modern font
% \usepackage{newtxtext,newtxmath}




% %--------------------------------------------------------------------------
% %         General Informations
% %--------------------------------------------------------------------------
\newcommand{\BigTitle}{
    \begin{center}
    华中科技大学微积分（B）考试合集
    \end{center}
    }

\newcommand{\LittleTitle}{
    那，边。
    }
\begin{document}
\frontmatter
\pagenumbering{Roman}

% %--------------------------------------------------------------------------
% %         First pages 
% %--------------------------------------------------------------------------
\newgeometry{top=8cm,bottom=.5in,left=2cm,right=2cm}
\subfile{files/0.0.0.titlepage}
\restoregeometry
\subfile{files/0.Preface}
\subfile{files/0.zommaire}
\subfile{files/0.1.changelogs}

% %--------------------------------------------------------------------------
% %         Core of the document 
% %--------------------------------------------------------------------------
% \subfile{files/mainpoints}
\mainmatter
\part{微积分（B）（上）期中考试}
\subfile{files/2025}
\subfile{files/2024}
\subfile{files/2023}
\subfile{files/2022}
\subfile{files/2021}
\subfile{files/2020}
\subfile{files/2019}
\subfile{files/2018}
\subfile{files/2017}
\subfile{files/2016}
\subfile{files/2015}
\subfile{files/2014}
\subfile{files/2013}
\subfile{files/2012}
\subfile{files/2011}
\part{微积分（B）（上）期末考试}
\subfile{files/2025f}
\subfile{files/2024f}
\subfile{files/2023f}
\subfile{files/2022f}
\subfile{files/2021f}
\subfile{files/2019f}
\subfile{files/2018f}
\subfile{files/2016f}
\part{微积分（B）（下）期中考试}
\subfile{files/2026x}
\subfile{files/2025x}
\subfile{files/2024x}
\subfile{files/2023x}
\subfile{files/2022x}
\subfile{files/2021x}
\subfile{files/2019x}
\subfile{files/2018x}
\subfile{files/2017x}
\part{微积分（B）（下）期末考试}
\subfile{files/2025xf}
\subfile{files/2024xf}
\subfile{files/2023xf}
\subfile{files/2022xf}
\subfile{files/2021xf}
\subfile{files/2019xf}
\subfile{files/2018xf}
\subfile{files/2017xf}
\subfile{files/2016xf}

% \part{Second part}
% \subfile{files/0.0.testchap}
% \subfile{files/0.0.testchap}
% \subfile{files/0.0.testchap}

% \part{Third part}
% \subfile{files/0.0.testchap}
% \subfile{files/0.0.testchap}
% \subfile{files/0.0.testchap}



% %--------------------------------------------------------------------------
% %         Bibliographie 
% %--------------------------------------------------------------------------
\nocite{*} % to cite evey things, else cite each on using : \cite{ifrs17}. 
% \printbibliography %to print bibliographie

\end{document}
